% ============================================================
%  Section 5: Conclusions
% ============================================================

This paper compared a slew-rate-constrained NMPC, with true-state access (Oracle) and with an augmented EKF/moving-horizon estimator (NMPC-MHE), against the two strongest robust PID tunings of an exothermic CSTR, AMIGO and Optimal ITAE with $M_s \le 1.6$, all acting through the same saturating, rate-limited valve with directional anti-windup on the PID side. Five criteria delimit when the computational overhead of nonlinear control pays off on this reactor:
\begin{enumerate}
    \item \textbf{Near the design point, an $M_s$-constrained PID is sufficient.} The NMPC improves the nominal servo IAE of Optimal ITAE by $2.3\%$ and the wide-range tracking IAE of AMIGO by $2.5\%$. The relevant baseline for a nonlinear controller is a robustly tuned PID with correct anti-windup.
    \item \textbf{Nonlinear prediction pays off after a change of the plant that the model represents.} After an abrupt loss of heat transfer to $70\%\,UA$, NMPC-Oracle reduces the IAE of the best PID (AMIGO) by $21.0\%$ and that of Optimal ITAE by $25.2\%$, with less than half of their valve travel; NMPC-MHE stays within $1.7\%$ of the Oracle. In wide-range tracking the Oracle needs about four times less valve travel than the PID loops.
    \item \textbf{Without a disturbance model, nonlinear prediction does not help regulation.} Under unmeasured feed shocks both NMPC variants settle with an offset of $+8.69$ and $+8.91$~K, while both PID loops return to the setpoint. Offset-free formulations with integrating disturbance states are a prerequisite for using NMPC as a regulator on this process.
    \item \textbf{Output feedback removes part of the advantage.} The Oracle is invariant to sensor noise; NMPC-MHE, the implementable variant, uses valve travel comparable to that of a PID loop once $\sigma \ge 0.1$~K, although it still converts less noise into thermal deviation ($\mathrm{IAE} = 3.33$ versus $4.16$~K$\cdot$min at $\sigma = 1.0$~K).
    \item \textbf{No controller crosses the thermodynamic boundary.} At $60\%\,UA$ all four controllers give practically the same error ($30.56$--$31.08$~K$\cdot$min) with the valve driven to its limit; the boundary is set by the heat-removal capacity of the jacket, not by the control law.
\end{enumerate}
The cost of these benefits is moderate: $31.46$ and $39.28$~ms per sample for NMPC-Oracle and NMPC-MHE against $0.09$~ms for a PID loop, well within the $1.2$~s sampling interval. The practical recommendation is therefore to adopt NMPC on such reactors when operation is dominated by changes of the plant, such as fouling or wide operating ranges, and to retain a robust PID, or to add an integrating disturbance model to the NMPC, when operation is dominated by unmeasured feed disturbances. Future work will add offset-free disturbance models to the NMPC and test the criteria on reactors with structural model mismatch.
